\chapter{Mynotes from 2026 class}
Some references used here:
\begin{itemize}
    \item Mumford's red book~\cite{mumford1988red}
    \item Humphreys Linear algebraic groups~\cite{humphreys1975linear}
\end{itemize}

\section{Lecture 1 and 2}
First some definitions and lemmas worth to remember.
\begin{itemize}
    \item (Integral extension). Let $R$ be a ring and $S\subset R$ a subring. We say that 
    $R$ is integrally dependent over $S$ if every $r\in R$ is integral over $S$, that is, there 
    exist a monic polynomial $f\in S[X]$ such that $f(r)=0$.
    \begin{lemma}\label{lemma : characterization of integral extensions}
        Let $S\subset R$ be rings. 
        A finite subset of elements $b_1,\dots, b_m\in R$ are all integral over $S$ if and only if 
        $S[b_1,\dots, b_m]$ is finitely generated (as a $S-$module) over $S$.
    \end{lemma}
    Proof: $(\Rightarrow)$ Let $b\in R$ be integral over $S$, i.e., there exist 
    \begin{equation}
     f =    X^d + a_1X^{d-1} +\cdots + a_d \in S[X]
    \end{equation}
    such that $f(b)=0$. Then, using long division, we see that $S[b]=S +\cdots + S b^{d-1}$ which is finitely generated. 
    Now, if $b_1,\dots, b_m$ are all integral over $S$, we can argue by induction that $S[b_1,\dots, b_{m-1}]$ is finitely generated 
    over $S$. However, since $b_m$ is integral over $S[b_1,\dots, b_{m-1}],\;S[b_1,\dots,b_m]$ is finitely generated 
    over $[b_1,\dots, b_{m-1}]$ and it must be finitely generated over $S$ aswell.\par  
    $(\Leftarrow)$ Suppose that $S[b_1,\dots, b_m] = S\beta_1+\cdots+S\beta_n$ for some $\beta_i\in R$. Then for any 
    $b\in R$ which fixes $S[b_1,\dots,b_m]$, we have $b\beta_i= \sum_j a_{ij}\beta_j$ for some $(a_{ij})\in\text{Mat}_n(S)$. 
    Therefore 
    \begin{equation}
      0 =  (bI-(a_{ij}))^*(bI-(a_{ij}))v = \det(bI-(a_{ij}))v
    \end{equation}
    where $v=(\beta_1,\dots,\beta_n)^T$. Therefore $\det(bI-(a_{ij}))\beta_i=0,\forall i$ and since 
    $1\in S\beta_1+\cdots+S\beta_n$, we see that $\det(bI-(a_{ij}))=0$. Since $\det(X-(a_{ij}))\in S[X]$ is monic, we 
    are done.\qed 
    \begin{corollary}\label{corollary : transitive property integral dependence}
        (Transitive property of integral dependence). Given rings $R \supset S\supset T$, if $R$ is integral over $S$ and 
        $S$ is integral over $T$, then $R$ is integral over $T$.\par 
        \textcolor{red}{Question: Do we need Noetherian hypothesis using this statement?}
    \end{corollary}
    Proof: Let $r\in R$, then by hypothesis there exist $f=X^d+s_1X^{d-1}+\cdots\in S[X]$ such that $f(r)=0$.
    Define $A=T[s_1,\dots,s_d]$, then $A[r]$ is a finitely generated $A-$module because $S\subset A$ and lemma~(\ref{lemma : characterization of integral extensions}). 
    Now, since $A$ is a finitely generated $T-$module by lemma~(\ref{lemma : characterization of integral extensions}),
    we see that $A[r]$ is a finitely generated $T-$module. Now, \textcolor{red}{assuming $T$ is Noetherian} we see that 
    $T[r]$ is finitely generated as a $T-$module.\qed

    \begin{remark}
        We can certainly say (using the previous proof) that if $r\in R$, then $S[r]$ is finitely generated 
        as a $T-$module.
    \end{remark}

    \begin{lemma}\label{lemma : extension field transcendence degree}
        Let $K\supset k$ ($k$ a perfect field) be a finitely generated extension field with transcendence degree $\partial_k(K)=d$, then 
        there exist generators $z_1,\dots,z_d,z_{d+1}$ such that 
        \begin{itemize}
            \item[(1)] $K=k(z_1,\dots,z_{d+1})$
            \item[(2)] $\{z_1,\dots,z_d\}$ are algebraicaly independent
            \item[(3)] $z_{d+1}$ is separable over $k(z_1,\dots, z_d)$ 
        \end{itemize} 
    \end{lemma}  
\end{itemize}

\begin{dang}
\begin{lemma}
    (Noether's normalization lemma - Mumford's red book). Let $R$ be a finitely generated algebra over $k$. Assume also that $R$ is an integral domain.
    If $R$ has transcende degree $n$ over $k$ (which means that $\partial_k\text{Frac}(R)=n$), then 
    there exist $x_1,\dots, x_n\in R$ algebraicaly independent and such that $R$ is integral 
    over $k[x_1,\dots,x_n]$.
\end{lemma}
\end{dang}
\TODO Study prove this version of the lemma using the statement made in class.

\begin{dang}
    \begin{lemma}\label{lemma : Noether normalization}
        (Noether's normalization lemma - Statement made in class). Let $R$ be a finitely generated 
        $k-$algebra. Then there exist $x_1,\dots,x_n\in R$ algebraically independent over $k$ such that 
        $R$ is integral over $k[x_1,\dots,x_n]$.
    \end{lemma}
\end{dang}
Proof (Nagata): Since $R$ is finitely generated as a $k$-algebra, there exist generators $y_1,\dots,y_m$ such that 
$R= k[y_1,\dots, y_m]$. Now suppose that $y_1,\dots, y_m$ are algebraically independent, then we can take $x_i = y_i$. 
If not, then there exist $p\in k[Y_1,\dots, Y_m]\setminus0$ such that $p(y_1,\dots,y_m)=0$. Write 
\begin{equation}
    p(Y_1,\dots,Y_m)=\sum_{\alpha\in\bbZ^m_\geq0}a_\alpha Y_1^{\alpha_1}\cdots Y_m^{\alpha_m}
\end{equation}
and define a new polynomial 
\begin{equation}
    \begin{split}
        q(Y_1,\dots,Y_m) &= p(Y_1,Y_2+Y_1^{r_2},\dots, Y_m+Y_1^{r_m})\\
        &=\sum_\alpha a_\alpha(Y_1^{\alpha_1+r_2\alpha_2+\cdots+r_m\alpha_m}+\cdots) 
    \end{split}
\end{equation}
for some choice of positive integers $r_2,\dots,r_m$. Take a positive number $b>\alpha_i,\;\forall \alpha_i$ such that 
$a_\alpha\neq0$. If we take $r_i = b^i$, then $\alpha_1+\alpha_2b+\cdots +\alpha_mb^{m-1}$ are all different numbers in base $b$. 
So, from different powers of the form $\alpha_1+r_2\alpha_2+\cdots+r_m\alpha_m$ we have a maximal one 
$d=\max\{\alpha_1+\cdots r_m\alpha_m:a_\alpha\neq0\}$ and 
\begin{equation}
    q(Y_1,\dots,Y_m) = aY_1^d + \text{smaller degree terms} \in k[Y_2,\dots, Y_m][Y_1],\;a\neq0.
\end{equation}
Now, if we let $z_i=y_i - y_1^{r_i}\in R,i=2,\dots,m$, we see that $q(y_1,z_2,\dots,z_m)=p(y_1,\dots,y_m)=0$ and 
$y_1$ is integral over $k[z_2,\dots, z_m]$. By induction on number of generators, there exist $x_1,\dots, x_n$
algebraically independent such that $k[z_2,\dots, z_m]$ is integral over $k[x_1,\dots, x_n]$. So, by lemma~\ref{lemma : characterization of integral extensions}
and its corollary~(\ref{lemma : characterization of integral extensions}), we see that $k[y_1,z_2,\dots,z_m]$ is integral over $k[x_1,\dots,x_m]$. Now, by definition, it's clear 
that $k[y_1,z_2,\dots,z_m]$ is a finitely generated $R-$module. So by the corollary~(\ref{corollary : transitive property integral dependence}) again, we 
see that $R$ is integral over $k[x_1,\dots,x_m]$.\qed 

\begin{dang}
    \begin{lemma}\label{lemma : R is a field, S is a field}
        Let $S\subset R$ be a ring extension. If $R$ is a field, which integral over $S$, then $S$ is a field. 
    \end{lemma}
\end{dang}
Proof: Let $s\in S\setminus0$ and $r\in R$ be its inverse. Since $R$ is integral over $S$, there exist
$f=X^d+a_1X^{d-1}+\cdots\in S[X]$ be such that $f(r)=0$. Therefore 
\begin{equation}
   0= s^{d-1}f(r) = r + a_1 + sa_2+\cdots + s^{d-1}a_0\Rightarrow r\in S.
\end{equation} 
\qed

\begin{remark}
    Let $L|K$ be a field extension. Then
    \begin{center}
        \begin{tikzcd}
            L|K \text{ finite extension}\arrow[r, bend left= 30] &  L|K \text{ algebraic extension}
            \arrow[l, bend left= 30, "\text{+ finitely generated extension}"]
        \end{tikzcd}
    \end{center}
\end{remark}

\begin{dang}
    \begin{theorem}\label{theorem : Zariski 1}
        (Zariski). Let $k$ be a field and $\fr{m}\subset R=k[X_1,\dots, X_m]$ be a maximal ideal. Then $L=R/\fr{m}$ is 
        a finite extension field of $k$.
    \end{theorem}
\end{dang}
Proof: By Noether normalization lemma~(\ref{lemma : Noether normalization}), there exist 
$a_1,\dots, a_n\in L$ algebraically independent such that $L$ is integral over $k[a_1,\dots,a_n]$. By lemma~(\ref{lemma : R is a field, S is a field}), 
we conclude that $k[a_1,\dots,a_n]$ is a field. Suppose that $n>0$, so $a_1$ has an inverse. This means that 
$a_1g(a_1,\dots, a_n) = 1$ for some $g\in k[X_1,\dots, X_n]$, which constradicts the fact that $a_1,\dots, a_n$ are 
alegrbaically independent. Then we conclude that $L$ is an integral (and then algebraic) over $k$. However, since 
$L$ is a finitely generated extension ($L=k(\overline{X}_1,\dots,\overline{X}_m)$) over $k$ we conclude that $L$ is in fact 
a finite field extension over $k$.\qed

\begin{dang}
    \begin{theorem}\label{theorem : Hilbert's zero theorem}
        Suppose that $k$ is an algebraically closed field.
        Let $I\subset R =k[X_1,\dots, X_m]$ be an ideal, then $V(I)=\varnothing\iff I = k[X_1,\dots, X_m]$.
    \end{theorem}
\end{dang}
Proof: If $I\neq R$, then there exist a maximal ideal $\fr{m}\supset I$. So $V(I)\supset V(\fr{m})$ and it's sufficient to 
prove that $V(\fr{m})\neq\varnothing$. By theorem~(\ref{theorem : Zariski 1}), we have that $L=R/\fr{m}$ is a finite 
$k-$extension, so we must have $L=k$ because $k$ is algebraically closed. So there exist $a_i\in k,i=1,\dots,m$ such that 
$X_i - a_i\in\fr{m}$ and $(X_1-a_1,\dots, X_m-a_m)\subset \fr{m}$. However $\varphi: R\to k:F\mapsto F(a_1,\dots,a_m)$ has kernel 
$\Ker(\varphi)=(X_1-a_1,\dots, X_m-a_m)$ and since $k$ is a field, $R/\Ker(\varphi)$ must be  a field and 
$\Ker(\varphi)$ is a maximal ideal. Therefore we conclude that $\fr{m}=(X_1-a_1,\dots, X_m-a_m)$. Finally
$V(X_1-a_1,\dots,X_m-a_m)=\{a_1,\dots,a_m\}\neq\varnothing$.\qed

\begin{remark}
    The hypothesis $k=\overline{k}$ in the previous theorem is crucial. For example if $k=\bbR$, we have 
    $V(x^2+1)=\varnothing$.
\end{remark}

\begin{corollary}\label{corollary : maximal ideals in k[x1,...,xn] ring characterization}
    From the proof of the previous theorem, we see that there exist a bijection 
    \begin{equation}
        \begin{split}
            k^n &\longleftrightarrow \left\{\begin{matrix}\text{maximal ideals}\\\fr{m}\subset R\end{matrix}\right\}\\
           (a_1,\dots,a_n)&\longmapsto (X_1-a_1,\dots, X_n-a_n)
        \end{split}
    \end{equation}
\end{corollary}


\begin{dang}
    \begin{theorem}\label{theorem : Hilbert Nullstelensatz}
        (Hilbert Nullstelensatz). Let $A\subset k[X_1,\dots, X_n]$ be an ideal, then 
        \begin{equation}
            IV(A) = \sqrt{A}.
        \end{equation}
    \end{theorem}
\end{dang}
Proof: The inclusion $\sqrt{A}\subset IV(A)$ is easy. Now, assume $g\in IV(A)$, where $A=(f_1,\dots, f_m)$. This 
means that we have the hypothesis:
\begin{equation}
    f_i(a_1,\dots, a_n)=0,\;\forall i=1,\dots,m\Rightarrow g(a_1,\dots, a_n)=0\; (*)
\end{equation}
Now, define the ideal 
\begin{equation}
    I = Ak[X_1,\dots, X_{n+1}] + (1- g(X_1,\dots, X_n)X_{n+1}).
\end{equation}
If $I$ is a proper ideal, then by corollary~(\ref{corollary : maximal ideals in k[x1,...,xn] ring characterization}), we have 
$I\subset (X_1-a_1,\dots, X_{n+1}-a_{n+1})$, which is the kernel of the evaluation map 
$k[X_1,\dots,X_{n+1}]\rightarrow k: f\mapsto f(a_1,\dots, a_{n+1})$. This means that 
\begin{equation}
    \left\{\begin{aligned}
        &f_i(a_1,\dots,a_n)=0,\;\forall i=1,\dots,m\\
        &g(a_1,\dots,a_n)a_{n+1}=1
    \end{aligned}\right.
\end{equation}
which contradicts our initial hypothesis $(*)$. Therefore, $I$ is not a proper ideal, so $1\in I$ and there exists 
$h_1,\dots, h_{m+1}\in k[X_1,\dots, X_{n+1}]$ such that 
\begin{equation}
    1 = \sum_{j=1}^mh_j(X_1,\dots, X_{n+1})f_j(X_1,\dots,X_n)  + h_{m+1}(X_1,\dots,X_{n+1})(1-g(X_1,\dots,X_n)X_{n+1}).
\end{equation}
Now, by projecting this equation onto $k[X_1,\dots, X_{n+1}]/(1-g X_{n+1})$, we get 
\begin{equation}
    g^\ell(X_1,\dots,X_n) = \sum_{j=1}^m h_j(X_1,\dots, X_{n+1})g^\ell  \times f_j(X_1,\dots,X_n)
\end{equation}
modulo $(1-g X_{n+1})$, where $H_j = h_j(X_1,\dots, X_{n+1})g^\ell\in k[X_1,\dots,X_n],\;\forall j$ if $\ell\gg0$. 
Since $\{\overline{X}_1,\dots,\overline{X}_{n+1}\}$ is algebraically free in $k[X_1,\dots, X_{n+1}]/(1-g X_{n+1})$, we 
conclude that, in fact, we have equality 
\begin{equation}
     g^\ell(X_1,\dots,X_n) = \sum_{j=1}^m H_j(X_1,\dots,X_n)  \times f_j(X_1,\dots,X_n)
\end{equation}
in $k[X_1,\dots,X_n]$.\qed

\begin{remark}
    We could have used the homomorphism 
    \begin{equation}
        k[X_1,\dots,X_{n+1}]\rightarrow k(X_1,\dots,X_n):X_{n+1}\mapsto g,\;X_i\mapsto X_i,\;i\neq n+1
    \end{equation}
    at the end of the previous proof.
\end{remark}

\section{Lecture 3}
\begin{dang}
    \begin{definition}
        We define an affine algebraic set $V\subset k^n$ to be a set of the form $V=V(S)$ for 
        some $S\subset k[\underline{X}]$.
    \end{definition}
\end{dang}

\begin{notation}
    Let $V$ be an affine algebraic set. We write 
    \begin{equation}
        D_V(g) = \{p\in V: g(p)\neq0\} = V\setminus V(g).
    \end{equation}
    These open sets are called ``principal open subsets''.
\end{notation}

\begin{dang}
    \begin{lemma}
        Any open subset $U\subset V$ of an affine algebraic set is quasicompact.
    \end{lemma}
\end{dang}
Proof: Consider a covering $U=\bigcup_{\lambda\in\Lambda}U_\lambda$.
Observe that $V\setminus U=V(g_1,\dots, g_m)$, so we have the equality $U=\bigcup_{i=1}^m D_V(g_i)$. In particular, we see 
that principal open subsets form a basis for the Zariski topology on $V$. Therefore, 
it's sufficient to prove that a principal open subset is quasicompact. So we assume $U=D_V(g)$. Furthermore, we have 
$U_\lambda = \bigcup_{i=1}^{n_\lambda}D_V(g_{\lambda,i})$ and 
\begin{equation}
    U=\bigcup_{\lambda\in\Lambda}\bigcup_{i=1}^{n_\lambda}D_V(g_{\lambda,i}).
\end{equation}
So, without lost of generalizaiton, we can assume $U=\bigcup_{\lambda\in\Lambda}D_V(g_\lambda)$ from beginning. Now 
\begin{equation}
    \begin{split}
        D_V(g) & = \bigcup_{\lambda\in\Lambda}D_V(g_\lambda)\iff\\
      V\cap V(g) &= V\cap V(\langle g_\lambda:\lambda\in\Lambda\rangle)\iff\\ 
      V\cap V(g) &= V\cap V(\langle g_\lambda:\lambda\in I\rangle)\;(I\subset \Lambda\;\text{ finite}).
    \end{split}
\end{equation}
In the last equation we used Hilbert basis theorem: The sequence
\begin{equation}
    \langle g_{\lambda_1} \rangle \subset  \langle g_{\lambda_1},g_{\lambda_2} \rangle \subset \cdots 
\end{equation}
must terminate at some point, because $k[\underline{X}]$ is Noetherian. Then we conclude that 
$D_V(g)=\bigcup_{\lambda\in I}D_V(g_\lambda)$.\qed

\begin{dang}
    \begin{definition}
    (Regular functions). Let $V\subset k^n$ be an affine algebraic set and $U\subset V$ be an open set. 
    A map $\varphi:U\to k$ is said to be \textit{regular} if for any $p\in U$, there exist $p\in U^\prime\subset U$, 
    polynomials $f,g\in k[\underline{X}]$ such that $g(q)\neq0,\;\forall q\in U^\prime$ and  
    $\varphi|_{U^\prime}=f/g|_{U^\prime}$.    
    \end{definition}
\end{dang}

\begin{dang}
    \begin{notation}
    Let $V$ be an affine algebraic set and let $U\subset V$ be an open set, we denote 
    $\mathcal{F}(U)=\{U\overset{\text{regular}}{\rightarrow}k\}$.
\end{notation}
\end{dang}


\begin{dang}
    \begin{lemma}\label{lemma: regular functions on principal open sets}
        Let $V\subset k^n$ be an affine algebraic set and $g\in k[X_1,\dots, X_n]$. Then the ser of regular functions 
        $\varphi: D_V(g)\to k$ is of the form $f/g^m$ for some $f\in k[X_1,\dots, X_n]$ and $m>0$. So 
        \begin{equation}
            \mathcal{F}(D_V(g)) \cong k[X_1,\dots,X_n]_g
        \end{equation}
        where $\mathcal{F}$ denotes the sheaf of regular functions on $V$.
    \end{lemma}
\end{dang}
Proof: First we have $D_V(g)=\bigcup_{i=1}^\ell D_V(g_i)$ and since $\varphi$ is regular, there exists 
$f_i, h_i\in k[\underline{X}],\;i=1,\dots,\ell$ such that 
\begin{equation}
    q\in D_V(g_i)\Rightarrow h_i(q)\neq 0\;(\text{so }q\in D_V(h_i)),\;\varphi\big|_{D_V(g_i)}=\frac{f_i}{h_i}\bigg|_{D_V(g_i)}.
\end{equation}
Now, we have two fundamental relations:
\begin{equation}
    \left\{\begin{aligned}
        &D_V(g_i)\subset D_V(h_i),\;\forall i=1,\dots,\ell\\
        &D_V(g)=\bigcup_{i=1}^\ell D_V(g_i)
    \end{aligned}\right.
\end{equation}
Using $V=V(F_1,\dots, F_m)$ and Hilbert Nullstelensatz~(\ref{theorem : Hilbert Nullstelensatz}), we get 
\begin{equation}
    \left\{\begin{aligned}
        &\sqrt{(F_1,\dots,F_m,g_i)}\subset \sqrt{(F_1,\dots,F_m,h_i)},\;\forall i=1,\dots,\ell\quad (*)\\
        &\sqrt{(F_1,\dots,F_m,g)}=\sqrt{(F_1,\dots,F_m,g_1,\dots,g_\ell)}\quad(**)
    \end{aligned}\right.
\end{equation}
From $(*)$, we deduce that $g_i^{n_i}=\sum_j A_jF_j+B_ih_i$ and restricting to $V$ we get 
$g_i^{n_i}=B_ih_i$. Now, since $\varphi\big|_{D_V(g_i)}=\frac{B_i f_i}{g_i^{n_i}}$ and 
$D_V(g_i)=D_V(g_i^{n_i})$, we can \textcolor{red}{change notation} $B_i f_i\to f_i$ and $g_i^{n_i}\to g_i$.\par 

Now we do a trick: since $g_jf_i - f_jg_i = 0$ on $D_V(g_i)\cap D_V(g_j)=D_V(g_ig_j)$, we get 
\begin{equation}
    g_ig_j(g_jf_i-f_jg_i) = 0\;\text{ on }V.
\end{equation}
Furthermore, we still have $\varphi\big|_{D_V(g_i)}=g_if_i / g_i^2$.
So we \textcolor{red}{change notation} again: $g_i^2\to g_i$ and $g_if_i\to f_i$, so we have 
\begin{equation}
    g_if_j - g_jf_i = 0,\;\forall i,j=1,\dots,\ell\quad (***).
\end{equation}
Now, by equation $(**)$, there exist $m>0$ such that 
\begin{equation}
    g^m = \sum_jA_jF_j + \sum_kB_kg_k 
\end{equation}
and we get on $V$:
\begin{equation}
    g^m = \sum_k B_kg_k.
\end{equation}
Now, by using equation $(***)$, we have 
\begin{equation}
    f_ig^m = \sum_kB_kf_ig_k = \sum_kB_kf_kg_i = (\sum_kB_kf_k)g_i.
\end{equation}
So if $q\in D_V(g_i)$, we have $\varphi(q) = \tfrac{f_i}{g_i}(q)=\tfrac{f}{g^m}(q)$, where we set $f=\sum_kB_kf_k$.\qed

\begin{dang}
    \begin{corollary}
        Since $V=D_V(1)$, we have $\mathcal{F}(V)=k[X_1,\dots,X_n]$. So the set of regular 
        functions is the set of polynomials on $V$.
    \end{corollary}
\end{dang}

\begin{dang}
    \begin{definition}
        A morphism between two affine algebraic sets $V_1\subset k^n$ and $V_2\subset k^m$ is a function 
        $\varphi:V_1\to V_2$ such that for any open set $U\subset V_2$, we have 
        \begin{equation}
            \varphi_U^{\#}:\mathcal{F}(U,k)\overset{-\circ\varphi}{\rightarrow}\mathcal{F}(\varphi^{-1}(U),k),\;
            \varphi_U^{\#}(\mathcal{O}_{V_1}(U))\subset \mathcal{O}_{V_2}(\varphi^{-1}(U)).
        \end{equation}
    \end{definition}
\end{dang}

Using more fancing words, the previous definition says that a continuous map $\varphi:V_1\to V_2$ is a morphism 
if the composition map $\varphi_U^{\#}:\mathcal{F}(U,k)\overset{-\circ\varphi}{\rightarrow}\mathcal{F}(\varphi^{-1}(U),k)$
induces a sheaf morphism 
\begin{equation}
    \varphi^\#(U):\mathcal{O}_{V_2}\to f_*\mathcal{O}_{V_1}.
\end{equation}
We use the notation $f_*\mathcal{O}_{V_1}$ for direct image sheaf, it's defined by 
$f_*\mathcal{O}_{V_1}(U)=\mathcal{O}_{V_1}(f^{-1}(U))$.


\section{Lecture 4}

\begin{dang}
\textbf{First exercise}. Let $U=k^2-\{0\}$, show that the set of regular functions $U\to k$ are polynomials.    
\end{dang}
\textbf{Solution}. Let $\varphi\in\mathcal{O}_U(U)$ be a regular function, then 
$\varphi|_{D_{k^2}(X)}=f/X^m$ and $\varphi|_{D_{k^2}(Y)}=g/Y^n$ and on the intersection $D(X)\cap D(Y)=D(XY)$, we have 
\begin{equation}
    Y^nf - X^mg = 0\quad(*).
\end{equation}
Since there exist an infinity number of points $p\in D(XY)$, we also have the equality $(*)$ in 
$k[X,Y]$. Moreover, from this equation, we see that all the terms of $g$ are of the form 
$X^kY^\ell,\;\ell\geq n$ and all terms of $f$ are of the form $X^kY^\ell,\;k\geq m$. Therefore, the quotient
$F=f/X^m=g/Y^n$ is a polynomial and since $U = D(X)\cup D(Y)$, we get $\varphi = F\in k[X,Y]$.

\begin{remark}
    We used the fact that $\mathcal{O}_{U}(U^\prime)=\mathcal{O}_{k^2}(U^\prime)$ for all 
    $U^\prime\subset U$ open, which follows by definition. 
    The we applied lemma~(\ref{lemma: regular functions on principal open sets}).
\end{remark}

\begin{dang}
    \begin{definition}
        A \underline{function space} is a topological space $X$ with a sheaf $\mathcal{O}_X$ such that 
        for any open set $U\subset X$, we have 
        \begin{itemize}
            \item[(1)] $\mathcal{O}_X(U)\overset{\text{subring}}{\subset} \{U\overset{\text{continuous}}{\to}k\}$
            \item[(2)] Let $f\in\mathcal{O}_X(U)$, then $f$ has an inverse in $\mathcal{O}_X(U)$ if and only if 
            $f(q)\neq0,\;\forall q\in U$.
        \end{itemize}
    \end{definition}
\end{dang}


\begin{dang}
\begin{lemma}
    An affine algebraic set $V\subset k^n$ equiped if the sheaf of regular functions is a function space.
\end{lemma}
\end{dang}
Proof: Let $U\subset X$ be an open set and $f\in\mathcal{O}_X(U)$. Then, by definition, there exist an open cover 
$U=\bigcup_\lambda U_\lambda$ such that $f|_{U_\lambda}=f_\lambda/g_\lambda$, where $g_\lambda(q)\neq0\;\forall q\in U_\lambda$. 
To prove that $f$ is continuous, it's sufficient to prove that $f^{-1}(a)\subset V$ is closed for all $a\in k$. 
So we prove that $(f|_{U_\lambda})^{-1}(a)$ is closed in $V$ for all $\lambda$, which can be verified by the 
equality
\begin{equation}
    V\cap V(f_\lambda - ag_\lambda) = (f|_{U_\lambda})^{-1}(a).
\end{equation}
The second item is clear.\qed 

\begin{dang}
    \begin{definition}
        Let $(X,\mathcal{O}_X)$ and $(Y,\mathcal{O}_Y)$ be two function spaces. We define a morphism $X\to Y$ to be a 
        a continuous map $f:X\to Y$ which induces a sheaf map $f^\#:\mathcal{O}_Y\to f_*\mathcal{O}_Y$ in the sence that 
        if $f^\#(U)(g)=g\circ f$ for $g\in\mathcal{O}_Y(U)$, then $f^\#(U)(\mathcal{O}_Y(U))\subset \mathcal{O}_X(f^{-1}(U))$.
    \end{definition}
\end{dang}

\begin{dang}
    \begin{definition}
        An \underline{affine algebraic variety} is a function space $(X,\mathcal{O}_X)$ isomorphic to 
        an affine algebraic set with Zariski topology and sheaf of regular functions.
    \end{definition}
\end{dang}


\begin{dang}
    \begin{definition}
        Let $(X,\mathcal{O}_X)$ be a function space, and $U\subset X$ be an open set with induced topology and 
        $\mathcal{O}_U=\mathcal{O}_X|_U$. Then $(U,\mathcal{O}_U)$ is a function space called 
        \underline{induced open subspace}.
    \end{definition}
\end{dang}

\begin{dang}
    \begin{definition}
        An \underline{abstract algebraic variety} is a function space $(X,\mathcal{O}_X)$ such that there exist an open covering  
        $X=\bigcup_\lambda U_\lambda$ with the property that for each $\lambda$, the induced 
        open subspace $(U_\lambda,\mathcal{O}_{U_\lambda})$ is an affine algebraic variety.
    \end{definition}
\end{dang}

\begin{dang}
    \begin{remark}
    In many texts we have the additional hypothesis of irredutibility in the definition of algebraic varieties.
\end{remark}
\end{dang}

\begin{dang}
    \begin{definition}
        An affine algebraic variety is said to be \underline{quasi-affine} if it's isomorphic, as a function space, 
        to an induced open subspace of an affine algebraic variety.
    \end{definition}
\end{dang}

\begin{dang}
    \begin{lemma}
        Let $(X,\mathcal{O}_X)$ be a function space and $Y\subset k^n$ be a quasi affine algebraic set
    \end{lemma}
\end{dang}


\begin{dang}
    \begin{lemma}\label{lemma: coordinate criteria to a map be a morphism of function spaces}
        Let $(X,\mathcal{O}_X)$ be a function space and $Y\subset k^n$ be a quasi-affine algebraic set. 
        Then $f:X\to Y$ is a morphism if and only $f=(g_1,\dots,g_n)$ with $g_i\in\mathcal{O}_X(X)$. Here, 
        we think $Y$ as an induced open subspace of an affine algebraic set.
    \end{lemma}
\end{dang}
Proof: $(\Rightarrow)$ The coordinate functions $x_i|_Y:Y\to k$ are regular functions, by definition. Now 
$g_i = x_i\circ f = f^\#(Y)(x_i)\in\mathcal{O}_X(f^{-1}(Y))=\mathcal{O}_X(X)$, because 
$x_i\in\mathcal{O}_Y(Y)=\{Y\overset{\text{regular}}{\to}k\}$.\par 
$(\Leftarrow)$ A closed set of $Y$ is of the form $V(F)\cap Y$ 
\textcolor{red}{(would not be of the form $V(F_1,\dots, F_m)$ in general ?)}. However
$f^{-1}(V(F)\cap Y)=\{p\in X: F(g_1,\dots, g_n)(p)=0\}$ which is closed in $X$ provided we show that 
$F(g_1,\dots, g_n)$ is continuous. However, since $\mathcal{O}_X(X)$ is a ring, we have $F(g_1,\dots, g_n)\in\mathcal{O}_X(X)$. 
Furthermore, if $U\subset Y$ is an open set and $\varphi:U\to k$ is regular, we need to show that 
$f^{\#}(U)(\varphi)\in\mathcal{O}_X(f^{-1}(U))$. Let $U^\prime\subset U$ be such that 
$\varphi|_{U^\prime}=F/G,\;G(q)\neq0\;\forall q\in U^\prime$. Then we have 
\begin{equation}
    f^\#(U^\prime)(\varphi)=\varphi\circ f|_{f^{-1}(U^\prime)} = 
    \frac{F(g_1,\dots, g_n)}{G(g_1,\dots, g_n)}\bigg|_{f^{-1}(U^\prime)}.
\end{equation}
Furthermore 
\begin{equation}
    G\circ f(q) = G(g_1,\dots, g_n)(q) \neq 0,\;\forall q\in f^{-1}(U^\prime).
\end{equation}
Since $G(g_1,\dots,g_n)|_{f^{-1}(U^\prime)}\in\mathcal{O}_X(f^{-1}(U^\prime))$, the previous equation 
shows that the inverse $1/G(g_1,\dots,g_n)$ also is an element of $\mathcal{O}_X(f^{-1}(U^\prime))$. Now, let 
$U=\bigcup_\lambda U_\lambda,\;\varphi|_{U_\lambda}=F_\lambda/G_\lambda$, then we see that 
\begin{equation}
    f^\#(U_\lambda)(\varphi|_{U_\lambda})\subset \mathcal{O}_X(f^{-1}(U_\lambda)).
\end{equation}
The glueing axiom shows that $f^\#(U)(\varphi)\in\mathcal{O}_X(f^{-1}(U))$.\qed


\begin{dang}
    \begin{lemma}
        Let $V\subset k^n$ be an affine algebraic set and $D_V(F)\subset V$ an affine algebraic set. Then 
        the induced open subspace 
        $(D_V(F), \mathcal{O}_{D_V(F)})$ is a quasi-affine algebraic variety, by definition. However, we have more, 
        it's also an affine algebraic variety.
    \end{lemma}
\end{dang}
Proof: Let $V=V(G_1,\dots, G_m)$ and define the map 
\begin{equation}
    \varphi: D_V(F)\rightarrow V(x_0F-1, G_1,\dots, G_m)\subset k^{n+1}:(a_1,\dots,a_n)\mapsto 
    (\frac{1}{F(a_1,\dots, a_n)},a_1,\dots, a_n)
\end{equation}
By the lemma~(\ref{lemma: coordinate criteria to a map be a morphism of function spaces}), we see that 
$\varphi$ is a morphism. Furthermore, it has a set theoretical inverse given by the projection 
\begin{equation}
    \psi:V(x_0F-1, G_1,\dots, G_m)\to D_V(F):(a_0,\dots,a_n)\mapsto (a_1,\dots, a_n).
\end{equation}
By the same lemma, this map is also a morphism. Since $f\psi = \mathds{1}_{D_V(x0F-1,\dots)}$ and 
$\psi f = \mathds{1}_{D_V(F)}$, we see that $\varphi$ is an isomorphism (the induced 
sheaf morphism would be an isomorphism by construction).\qed

\begin{dang}
    \begin{remark}
        Let $(X,\mathcal{O}_X)$ be an affine algebraic variety and 
        $\varphi:(V,\mathcal{O}_V)\to (X,\mathcal{O_X})$ an isomorphism, where $v\subset k^n$ is an algebraic set. 
        Then this isomorphism, induces another isomorphism
        \begin{equation}
            (D_V(F),\mathcal{O}_{D_V(F)}) \to (X_f,\mathcal{O}_{X_f})
        \end{equation}
        where $F=f\circ\varphi$ and $X_f = \varphi(D_V(F))=\{p\in X: f(p)\neq0\}$.
    \end{remark}
\end{dang}

\begin{dang}
    \begin{corollary}
        A quasi-affine algebraic variety is an abstract algebraic variety.
    \end{corollary}
\end{dang}
Proof: A quasi-affine algebraic variety has a topological basis consisting of principal open subsets, which 
are affine algebraic varieties.